\section{符号说明}
\label{sec:符号说明}

本书采用如下字型约定：\textbf{向量一律用粗体斜体（$\boldsymbol{}$）表示}，例如位置向量 $\boldsymbol{r}$、切向量 $\boldsymbol{r}'(t)$、$\boldsymbol{r}_u$；\textbf{标量、分量与矩阵元素不加粗}，例如 $x(t)$、$E$、$F$、$G$、$\kappa$、$\tau$。\textbf{单位向量（单位向量组）用 $\mathbf{}$（正粗体）表示}，例如 Frenet 标架 $\{\mathbf{T},\mathbf{N},\mathbf{B}\}$、曲面单位法向量 $\mathbf{n}$、活动标架 $\mathbf{e}_i$。空间、点、曲线、曲面等对象用正常体。

此处列出本笔记所用的主要符号（可参照《符号约定与工作要求.md》）：

\begin{table}[H]
\centering
\caption{符号说明}
\label{tab:符号说明}
\begin{tabular}{ll}
\toprule
符号 & 含义 \\
\midrule
$\mathbb{E}^3$ & 三维欧氏空间 \\
$\boldsymbol{r}(t)$ & 曲线的位置向量函数 \\
$\boldsymbol{r}(u,v)$ & 曲面的位置向量函数 \\
$\boldsymbol{r}_u,\boldsymbol{r}_v$ & 曲面的偏导向量（切向量） \\
$\mathbf{n}$ & 曲面的单位法向量 \\
$s$ / $ds$ & 弧长参数 / 弧长元素 \\
$\kappa$ / $\tau$ & 曲率 / 挠率 \\
$\rho$ & 曲率半径（$\rho=1/\kappa$） \\
$\mathbf{T},\mathbf{N},\mathbf{B}$ & 单位切向量、主法向量、次法向量（Frenet 标架） \\
$E,F,G$ & 第一基本形式系数 \\
$L,M,N$ & 第二基本形式系数 \\
$\mathrm{I}$ / $\mathrm{II}$ & 第一 / 第二基本形式 \\
$\kappa_1,\kappa_2$ & 主曲率 \\
$K$ / $H$ & 高斯曲率 / 平均曲率 \\
$\Gamma_{ij}^{k}$ & 第二类 Christoffel 符号 \\
$\omega^{i}$ / $\omega_{i}^{j}$ & 平移形式 / 连接形式 \\
$\wedge$ / $\mathrm{d}$ & 外积 / 外微分算子 \\
\bottomrule
\end{tabular}
\end{table}

符号表中 $E,F,G$ 与 $L,M,N$ 分别为第一、第二基本形式的系数（标准记号）；$N$ 作为系数使用时与主法向量 $\mathbf{N}$ 以正粗体区分。
